\section{Experimental Setup}

This chapter presents the quantitative evaluation of the proposed 
hybrid FAIRness assessment framework. The results assess the 
reliability of the framework and compare its outputs against 
manually annotated ground truth and an existing baseline approach.

\subsection{Dataset}

The full dataset comprises 188 research software repositories
collected from GitHub across 10 research domains, as described in
Chapter~3. In addition, a manually annotated corpus of 30 repositories
was constructed to provide ground truth for quantitative evaluation.
The annotated corpus was selected to provide diversity across research
domains and programming languages, including repositories implemented
in Python, C++, Java, MATLAB, Fortran, Julia, Go, Rust, and Groovy.

Ten of the 30 annotated repositories were initially used during a
pilot phase to refine the annotation guidelines and assessment
procedure. These repositories were retained in the final annotated
corpus. To assess the effect of this original pilot overlap, the
remaining 20 repositories were additionally analysed as a
pilot-exclusion subset. For the LLM-based dimensions, these
repositories had not been used during pilot-phase refinement and
therefore provide a pilot-exclusion sensitivity check of LLM--human
agreement that is unaffected by the original pilot-phase refinement.
For the rule-based indicators, the same subset is used
only as an internal consistency check because some rule-based
annotations were corrected during final scorer validation.

\subsection{Evaluation Metrics}

The following metrics are used throughout this chapter:

\begin{description}
\item[Exact Match Accuracy] The proportion of repositories for 
which the automated score exactly matches the manual annotation. 
Reported as a percentage.

\item[Mean Absolute Error (MAE)] The average absolute difference 
between the automated assessment and the corresponding manual 
annotation, computed per dimension across all 30 ground truth 
repositories. Applied to the ordinal LLM-based dimensions (D1--D6).

\item[Cohen's Unweighted Kappa] Used for binary rule-based 
indicators (F1--R2), where any disagreement is treated equally 
regardless of magnitude.

\item[Cohen's Linearly Weighted Kappa] Used for ordinal AI-based 
dimensions (D1--D6), where partial agreement is given proportional 
credit. A disagreement of 1 point is penalised less than a 
disagreement of 2 points.
\end{description}

Kappa values are interpreted following the scale proposed by
Landis and Koch~\cite{landis1977kappa}: values below 0 indicate poor
agreement, 0.00--0.20 slight agreement, 0.21--0.40 fair agreement,
0.41--0.60 moderate agreement, 0.61--0.80 substantial agreement,
and 0.81--1.00 almost perfect agreement.

\subsection{Models and Tools}

The LLM-based component was evaluated using two model configurations.
The primary configuration uses \texttt{llama3.1:8b}, executed locally
through the Ollama HTTP API at temperature 0. The model was selected
because it can be executed locally without external API calls,
supporting reproducible assessment within the local experimental
environment. Structured JSON output is enforced through a schema that
constrains scores to the permitted values (0, 1, or 2).

To examine whether the observed LLM assessment behaviour was specific
to this model configuration, a second evaluation was conducted using
\texttt{Qwen3.6-35B-A3B-MLX-8bit}. The model was accessed through a
university-hosted OpenAI-compatible chat completion API and was applied
to the same 30-repository annotated corpus. The same D1--D6 dimensions,
scoring rubric, README-only evidence policy, 12,000-character input
limit, and temperature setting were retained. Backend-specific output
handling differed as described in Chapters~3 and~4.

The rule-based component operates deterministically on locally stored
repository snapshots. The howfairis tool is used as a software-oriented
baseline comparison. F-UJI, an automated FAIR assessment tool designed
primarily for research data objects, is additionally used in a pilot
experiment to investigate the applicability of an external,
data-oriented FAIR assessment baseline to research software
repositories.

\subsection{Experimental Environment}

The evaluation pipeline was implemented in Python. The primary Llama
configuration was executed locally on Windows using Ollama, with the
\texttt{llama3.1:8b} model (4.9~GB). The second model configuration,
\texttt{Qwen3.6-35B-A3B-MLX-8bit}, was accessed through the
university-hosted OpenAI-compatible API. All assessment scripts, input
data, parsed outputs, and derived evaluation results are stored in the
project repository to support reproducibility. Model-specific result
files are kept separate so that evaluations can be reproduced and
analysed independently.


\subsection{Evaluation Procedure}

The evaluation was conducted in the following order:

\begin{enumerate}

\item The thesis author manually annotated 30 ground truth
repositories against the finalised annotation guidelines (v1.4).

\item A second annotator independently annotated the same 30
repositories using the same guidelines.

\item Inter-annotator agreement was computed between the two
human annotators to validate the reliability of the ground truth.

\item The Llama-based scorer was applied to the 30 annotated
repositories and its outputs were compared against both human
annotators.

\item The Qwen-based scorer was independently applied to the same
30 repositories using the same documentation dimensions and scoring
rubric, and its outputs were likewise compared against both human
annotators.

\item The rule-based scorer was applied to all 188 repositories
to produce a complete set of structural FAIR indicator scores.

\item Hybrid FAIRness profiles were generated by combining the
shared rule-based results with the corresponding LLM results for
the 30 ground truth repositories. Separate profile sets were retained
for the two evaluated LLM configurations.

\item The howfairis baseline was executed on the 30 ground truth
repositories and compared against the proposed framework.

\item F-UJI was applied directly to a subset of five ground truth
repositories to investigate its applicability as an external
data-oriented FAIRness assessment baseline for research software.

\item As a sensitivity analysis, LLM--human agreement was recomputed
for both model configurations on the 20-repository pilot-exclusion
subset after excluding the 10 repositories used during pilot-phase
refinement. Rule-based performance was also recomputed on this subset
as an internal consistency check.

\end{enumerate}

\section{Ground Truth Validation}

\subsection{Annotation Process}

Manual annotations were created by the thesis author for all 30 
ground truth repositories following the finalised annotation 
guidelines (v1.4), as described in Chapter~3. Each repository was 
annotated by examining the repository root file structure and README 
content. Rule-based indicators (F1--R2) were assigned binary scores 
according to the annotation guidelines based on evidence obtained 
from the repository structure and README content. AI-based dimensions 
(D1--D6) were scored on a 0--2 ordinal scale based on the clarity 
and completeness of the README documentation.

To assess annotation reliability, a second annotator independently 
annotated the same 30 repositories using identical guidelines and 
pre-fetched repository data, without access to the primary 
annotations.

\subsection{Inter-Annotator Agreement}

Inter-annotator agreement between the two annotators was quantified 
using Cohen's unweighted Kappa for binary rule-based indicators and 
Cohen's linearly weighted Kappa for ordinal AI-based dimensions. 
Table~\ref{tab:iaa} presents the results.

\begin{table}[h]
\centering
\caption{Inter-annotator agreement for rule-based indicators and 
AI-based dimensions (n=30).}
\label{tab:iaa}
\begin{tabular}{llrrr}
\hline
\textbf{Indicator/Dimension} & \textbf{Type} & \textbf{$\kappa$} & 
\textbf{Exact} & \textbf{Interpretation} \\
\hline
F1 --- Persistent identifier & Binary & 0.710 & 90.0\% & Substantial \\
F2 --- Structured metadata   & Binary & 1.000 & 100.0\% & Almost perfect \\
A1 --- License               & Binary & undef. & 100.0\% & Single class \\
I1 --- Dependency file       & Binary & 0.634 & 93.3\% & Substantial \\
R1 --- README present        & Binary & undef. & 100.0\% & Single class \\
R2 --- Citation information  & Binary & 0.658 & 83.3\% & Substantial \\
\hline
\textbf{Rule macro-avg}      & & \textbf{0.750} & & \\
\hline
D1 --- Installation          & Ordinal & 0.539 & 70.0\% & Moderate \\
D2 --- Usage                 & Ordinal & 0.747 & 83.3\% & Substantial \\
D3 --- Purpose               & Ordinal & 0.789 & 96.7\% & Substantial \\
D4 --- Dependencies          & Ordinal & 0.743 & 80.0\% & Substantial \\
D5 --- Examples              & Ordinal & 0.585 & 70.0\% & Moderate \\
D6 --- Contribution          & Ordinal & 0.849 & 90.0\% & Almost perfect \\
\hline
\textbf{AI macro-avg}        & & \textbf{0.709} & & \\
\hline
\end{tabular}
\end{table}

Overall, the agreement results indicate that both the binary 
rule-based indicators and the ordinal documentation quality 
dimensions can be annotated consistently using the proposed 
guidelines.

The AI-based dimensions achieved a macro-average weighted Kappa of 
0.709, indicating substantial overall agreement between annotators. 
The strongest agreement was observed for D6 (contribution 
information, $\kappa$=0.849) and D3 (purpose description, 
$\kappa$=0.789). The weakest agreement was observed for D1 
(installation clarity, $\kappa$=0.539) and D5 (examples and 
tutorials, $\kappa$=0.585), suggesting greater ambiguity when 
assessing partially documented functionality or documentation 
provided through external links.

For the rule-based indicators, F2 achieved perfect agreement, while 
F1, I1, and R2 showed substantial agreement. Two indicators are 
mathematically undefined: R1 (README present), because all 30 
repositories have a README file, and A1 (license), because both 
annotators agreed on the presence of a license for all 30 
repositories following a correction described below. Both indicators 
demonstrate consistent annotation in practice, as reflected by their 
100\% exact agreement.

The rule-based macro-average of 0.750 excludes the two undefined 
values (A1 and R1).

I1 (dependency file present) initially showed disagreement on seven 
repositories, predominantly projects using non-Python language 
ecosystems (C++, Java, Go, Rust, Groovy, Julia), following the 
pattern that the primary annotator assigned a score of 0 while the 
second annotator assigned a score of 1. Investigation traced this to 
an incomplete file-detection schema used during ground-truth 
collection: the collector script for the 30 ground-truth repositories 
predated the expanded dependency-file list (including 
\texttt{go.mod}, \texttt{Cargo.toml}, \texttt{pom.xml}, 
\texttt{build.gradle}, and \texttt{Project.toml}) later adopted for 
the full 188-repository collection, and consequently could not 
detect language-specific dependency files for non-Python projects. 
Five of the seven affected repositories were confirmed via the 
corrected file inventory to genuinely contain a dependency 
specification file, and the corresponding primary annotations were 
corrected accordingly. Following this correction, I1 agreement 
improved from $\kappa$=0.286 (fair) to $\kappa$=0.634 (substantial).

A1 (license) was also corrected during validation: one repository 
(\texttt{rust-ml/linfa}) had been annotated as lacking a license, 
despite both the GitHub repository metadata and the README text 
explicitly identifying an Apache-2.0 license. This annotation was 
corrected, resulting in full agreement between annotators for A1.

Overall, these inter-annotator agreement results indicate a level of
annotation consistency that supports using the manual annotations as
the reference ground truth for the subsequent evaluation, while the
dimension-level differences are considered when interpreting the
LLM--human agreement results.

\section{Rule-Based Assessment Results}

The rule-based scorer was applied to all 188 repositories in the 
full dataset, producing binary scores for each of the six FAIR 
indicators (F1--R2). This section presents the distribution of 
rule-based assessment results across the full dataset and 
evaluates the rule-based component against the manually annotated 
ground truth repositories.

\subsection{Full Dataset Results}

Table~\ref{tab:rb_full} summarises the frequency of positive scores 
for each indicator across all 188 repositories.

\begin{table}[h]
\centering
\caption{Rule-based indicator frequencies across the full dataset 
(n=188).}
\label{tab:rb_full}
\begin{tabular}{llrr}
\hline
\textbf{Indicator} & \textbf{Description} & \textbf{Positive} & 
\textbf{Percentage} \\
\hline
F1 & Persistent identifier & 95 & 51\% \\
F2 & Software metadata     & 58 & 31\% \\
A1 & License               & 182 & 97\% \\
I1 & Dependency file       & 123 & 65\% \\
R1 & README present        & 188 & 100\% \\
R2 & Citation information  & 79 & 42\% \\
\hline
\end{tabular}
\end{table}

All 188 repositories contain a README file (R1=100\%), and 97\% 
have a license file (A1=97\%). In contrast, dedicated software 
metadata files such as \texttt{CITATION.cff} or 
\texttt{codemeta.json} are present in only 31\% of repositories 
(F2), and persistent identifiers such as DOIs or Zenodo badges 
are present in only 51\% (F1). Citation information is available 
in 42\% of repositories (R2), suggesting that fewer than half of 
the evaluated repositories provide explicit guidance on how to 
cite the software. Overall, the results indicate that basic 
repository documentation practices, such as providing README and 
license files, are widespread, whereas richer FAIR-related metadata 
and citation support remain considerably less common.

The mean rule-based score across the full dataset was 3.86 out of 
a maximum possible score of 6. The score distribution is shown in 
Table~\ref{tab:rb_dist}.

\begin{table}[h]
\centering
\caption{Distribution of rule-based total scores across the full 
dataset (n=188).}
\label{tab:rb_dist}
\begin{tabular}{rrr}
\hline
\textbf{Score} & \textbf{Count} & \textbf{Percentage} \\
\hline
6/6 & 15  & 8\%  \\
5/6 & 43  & 23\% \\
4/6 & 61  & 32\% \\
3/6 & 42  & 22\% \\
2/6 & 23  & 12\% \\
1/6 & 4   & 2\%  \\
\hline
\end{tabular}
\end{table}

The score distribution spans the full range from 1 to 6, indicating
substantial variability in the structural FAIR-related indicators
assessed across the dataset. Only 15 repositories (8\%) satisfy all
six rule-based indicators, while 27 repositories (14\%) satisfy two
or fewer, indicating gaps in the structural evidence captured by the
assessment.


\subsection{Validation Against Ground Truth}

The rule-based scorer was validated against the final manual
annotations for the 30 ground truth repositories. During validation,
differences between automated predictions and manual annotations
were inspected individually and resolved according to the finalised
annotation guidelines (v1.4).

This process identified both implementation-level detection gaps and
manual annotation inconsistencies, particularly affecting the citation
(R2), persistent identifier (F1), and license (A1) indicators. The
rule-based detection logic was refined to recognise additional valid
citation and persistent identifier patterns, including
\texttt{CITATION.cff} files, generic DOI representations, and explicit
citation wording. Manual annotations that were inconsistent with the
final rubric were corrected and documented in the corresponding
repository annotation files.

After these refinements, the final rule-based scorer achieved exact
agreement with the final ground truth annotations across all 30
repositories. Precision, recall, and F1-score were 1.000 for each of
the six binary indicators (F1--R2), as shown in
Table~\ref{tab:rule_prf}.

\begin{table}[h]
\centering
\caption{Rule-based validation against manual ground truth (n=30).}
\label{tab:rule_prf}
\begin{tabular}{lrrr}
\hline
\textbf{Indicator} & \textbf{Precision} & \textbf{Recall} &
\textbf{F1-score} \\
\hline
F1 --- Persistent identifier & 1.000 & 1.000 & 1.000 \\
F2 --- Software metadata     & 1.000 & 1.000 & 1.000 \\
A1 --- License               & 1.000 & 1.000 & 1.000 \\
I1 --- Dependency file       & 1.000 & 1.000 & 1.000 \\
R1 --- README present        & 1.000 & 1.000 & 1.000 \\
R2 --- Citation information  & 1.000 & 1.000 & 1.000 \\
\hline
\end{tabular}
\end{table}

The final results demonstrate that the rule-based implementation
faithfully reproduces the operational definitions encoded in the
finalised manual annotation rubric for the evaluated ground-truth
corpus. Because scorer refinements and corrections to some rule-based
annotations were made during final validation, the resulting perfect
agreement should not be interpreted as an unbiased estimate of
performance on unseen repositories. The additional pilot-exclusion
analysis in Section~\ref{sec:pilot_exclusion_sensitivity} therefore
serves as an internal consistency check for the rule-based component 
rather than as independent validation.
\section{LLM-Based Assessment Results}
Two model configurations, \texttt{llama3.1:8b} and
\texttt{Qwen3.6-35B-A3B-MLX-8bit}, were evaluated on the same
30-repository ground-truth corpus. Both configurations assessed
the same six documentation dimensions (D1--D6) using the same scoring
rubric and README-only evidence policy.

The following analysis compares each model independently against both
human annotators. The purpose of the second model evaluation is not to
establish a general ranking between models, but to examine whether the
observed reliability and failure patterns are specific to the primary
Llama configuration or recur under a different LLM configuration.

\subsection{Agreement with the Primary Annotator}

Table~\ref{tab:llm_a1} presents the agreement of both LLM
configurations with the primary annotator. Weighted Cohen's Kappa,
exact agreement, and mean absolute error are reported independently
for each documentation dimension.

\begin{table}[htbp]
\centering
\caption{Agreement of the two LLM configurations with the primary
annotator (n=30).}
\label{tab:llm_a1}
\small
\begin{tabular}{lrrrrrr}
\hline
& \multicolumn{3}{c}{\textbf{Llama 3.1 8B}}
& \multicolumn{3}{c}{\textbf{Qwen3.6-35B-A3B-MLX-8bit}} \\
\cline{2-4}\cline{5-7}
\textbf{Dimension}
& \textbf{$\kappa$} & \textbf{Exact} & \textbf{MAE}
& \textbf{$\kappa$} & \textbf{Exact} & \textbf{MAE} \\
\hline
D1 --- Installation
& 0.605 & 70.0\% & 0.333
& 0.799 & 86.7\% & 0.133 \\
D2 --- Usage
& 0.274 & 46.7\% & 0.567
& 0.516 & 70.0\% & 0.300 \\
D3 --- Purpose
& $-$0.071 & 86.7\% & 0.133
& 0.789 & 96.7\% & 0.033 \\
D4 --- Dependencies
& 0.480 & 60.0\% & 0.433
& 0.487 & 56.7\% & 0.433 \\
D5 --- Examples
& 0.192 & 36.7\% & 0.700
& 0.320 & 43.3\% & 0.567 \\
D6 --- Contribution
& 0.494 & 63.3\% & 0.400
& 0.759 & 83.3\% & 0.167 \\
\hline
\textbf{Macro-avg. $\kappa$}
& \textbf{0.329} & & 
& \textbf{0.612} & & \\
\hline
\end{tabular}
\end{table}

The results differ substantially between the two model
configurations. Llama achieved a macro-average weighted Kappa of
0.329 against the primary annotator, whereas Qwen achieved 0.612.
The latter therefore showed substantially stronger agreement with
the manual annotations in this experiment.

For Qwen, the strongest agreement was observed for installation
(D1, $\kappa$=0.799), purpose (D3, $\kappa$=0.789), and contribution
information (D6, $\kappa$=0.759). Usage documentation achieved
moderate agreement (D2, $\kappa$=0.516), while dependencies
(D4, $\kappa$=0.487) and examples (D5, $\kappa$=0.320) remained
more difficult.

The contrast for D3 is particularly pronounced. Llama achieved a
weighted Kappa of $-$0.071 despite 86.7\% exact agreement, whereas
Qwen achieved $\kappa$=0.789 with 96.7\% exact agreement. The
negative Llama value occurs in a dimension with strongly imbalanced
human scores and should therefore be interpreted together with its
high exact agreement. At the same time, the substantially stronger
Qwen result shows that the weak Kappa observed for Llama is not a
consistent property of LLM-based purpose assessment across the two
evaluated configurations.

Overall, these results indicate that agreement with human assessment
is both model-dependent and dimension-dependent. The second model
substantially improves agreement for several dimensions, but the
remaining variation across D1--D6 shows that stronger overall
performance does not remove dimension-specific assessment
difficulties.

\subsection{Agreement with the Second Annotator}

To examine whether the model--human agreement patterns depend on the
choice of human reference, both model configurations were also
compared with the second annotator. Table~\ref{tab:llm_a2} presents
the results.

\begin{table}[htbp]
\centering
\caption{Agreement of the two LLM configurations with the second
annotator (n=30).}
\label{tab:llm_a2}
\small
\begin{tabular}{lrrrr}
\hline
\textbf{Dimension}
& \textbf{Llama $\kappa$}
& \textbf{Llama Exact}
& \textbf{Qwen $\kappa$}
& \textbf{Qwen Exact} \\
\hline
D1 --- Installation & 0.480 & 70.0\% & 0.429 & 63.3\% \\
D2 --- Usage        & 0.276 & 50.0\% & 0.488 & 66.7\% \\
D3 --- Purpose      & $-$0.056 & 86.7\% & 1.000 & 100.0\% \\
D4 --- Dependencies & 0.338 & 53.3\% & 0.481 & 60.0\% \\
D5 --- Examples     & 0.156 & 30.0\% & 0.133 & 13.3\% \\
D6 --- Contribution & 0.454 & 60.0\% & 0.809 & 86.7\% \\
\hline
\textbf{Macro-avg. $\kappa$}
& \textbf{0.275} & & \textbf{0.557} & \\
\hline
\end{tabular}
\end{table}

The comparison with the second annotator shows the same broad
model-level pattern. Llama achieved a macro-average weighted Kappa
of 0.275, whereas Qwen achieved 0.557. Qwen again showed particularly
strong agreement for purpose (D3, $\kappa$=1.000) and contribution
information (D6, $\kappa$=0.809).

However, the second-annotator comparison also highlights the
sensitivity of dimension-level results to the human reference.
Most notably, Qwen agreement for examples (D5) decreased to
$\kappa$=0.133, compared with 0.320 against the primary annotator.
This is consistent with the lower human--human agreement previously
observed for D5 and reinforces that some documentation dimensions
contain greater interpretive variation than others.

Taken together, the two annotator comparisons show that Qwen achieved
higher macro-average agreement than Llama against both human
references, while the magnitude of agreement remained dependent on
both the assessed dimension and the human reference annotation.

\subsection{Pilot-Exclusion Sensitivity Analysis}
\label{sec:pilot_exclusion_sensitivity}

Because 10 of the 30 annotated repositories had been used during
pilot-phase refinement, final-system results were additionally
examined after excluding these repositories. The resulting
20-repository subset provides a pilot-exclusion sensitivity check for
the LLM-based dimensions, which were not modified during this
validation process. For the rule-based component, however, the
analysis should be interpreted more cautiously: several manual
rule-based annotations within this subset were corrected during
final scorer validation after inconsistencies in the annotation
evidence and detection schema were identified. The rule-based
20-repository result is therefore reported as an internal consistency
check rather than as independent validation.

For the rule-based component, precision, recall, and F1-score were
1.000 for all six indicators after exclusion of the 10 original
pilot repositories. Because some rule-based annotations in this
subset were corrected during final validation, this result confirms
consistency between the final scorer and finalised annotation rubric
but should not be interpreted as an unbiased estimate of performance
on unseen repositories.

For the LLM-based component, Table~\ref{tab:llm_pilot_exclusion}
compares agreement with the primary annotator on the full annotated
corpus and the 20-repository pilot-exclusion subset for both model
configurations.

\begin{table}[htbp]
\centering
\caption{LLM agreement with the primary annotator on the full
annotated corpus and the 20-repository pilot-exclusion subset.}
\label{tab:llm_pilot_exclusion}
\small
\begin{tabular}{lrrrr}
\hline
\textbf{Dimension}
& \textbf{Llama (n=30)}
& \textbf{Llama (n=20)}
& \textbf{Qwen (n=30)}
& \textbf{Qwen (n=20)} \\
\hline
D1 --- Installation
& 0.605 & 0.659 & 0.799 & 0.787 \\
D2 --- Usage
& 0.274 & 0.474 & 0.516 & 0.531 \\
D3 --- Purpose
& $-$0.071 & $-$0.071 & 0.789 & 0.783 \\
D4 --- Dependencies
& 0.480 & 0.337 & 0.487 & 0.405 \\
D5 --- Examples
& 0.192 & 0.286 & 0.320 & 0.367 \\
D6 --- Contribution
& 0.494 & 0.556 & 0.759 & 0.815 \\
\hline
\textbf{Macro-avg.}
& \textbf{0.329}
& \textbf{0.373}
& \textbf{0.612}
& \textbf{0.615} \\
\hline
\end{tabular}
\end{table}

For Llama, excluding the 10 pilot repositories did not reduce
agreement with the primary annotator. The macro-average weighted
Kappa increased from 0.329 on the complete 30-repository corpus to
0.373 on the 20-repository pilot-exclusion subset. Agreement increased
for D1, D2, D5, and D6, remained unchanged for D3, and decreased for
D4.

The Qwen results were even more stable across the two populations.
Its macro-average weighted Kappa was 0.612 on the complete annotated
corpus and 0.615 after exclusion of the 10 pilot repositories.
Dimension-level agreement also remained broadly similar: D1 and D3
decreased only slightly, D2 and D5 increased slightly, D6 increased
from 0.759 to 0.815, and D4 decreased from 0.487 to 0.405.

Since neither LLM scorer nor the D1--D6 annotations in the
20-repository subset were modified as part of the rule-based
validation corrections, these results provide a sensitivity check
for the effect of the original pilot overlap on LLM--human agreement.
For both model configurations, exclusion of the pilot repositories
did not weaken the overall agreement result. In particular, the
near-identical Qwen macro-average values on the full corpus and the
pilot-exclusion subset indicate that its observed agreement was not
primarily driven by the 10 repositories involved in the original
pilot phase.

The rule-based results likewise remain internally consistent after
pilot exclusion. However, their perfect performance should not be
interpreted as independent generalisation evidence because some
rule-based annotations within this subset were corrected during final
scorer validation.

\subsection{Systematic Scoring Patterns}
\label{sec:llm_systematic_patterns}

Analysis of the direction of disagreements reveals that both model
configurations tended to assign lower scores than the primary human
annotator, although the dimensions in which this behaviour was most
pronounced differed between models. Table~\ref{tab:llm_dir}
summarises the direction of disagreements for both configurations.

\begin{table}[htbp]
\centering
\caption{Direction of LLM scoring disagreements relative to the
primary annotator (n=30). Lower indicates model score $<$ human score;
higher indicates model score $>$ human score.}
\label{tab:llm_dir}
\small
\begin{tabular}{lrrrr}
\hline
\textbf{Dimension}
& \textbf{Llama lower}
& \textbf{Llama higher}
& \textbf{Qwen lower}
& \textbf{Qwen higher} \\
\hline
D1 --- Installation  & 8  & 1 & 2  & 2 \\
D2 --- Usage         & 14 & 2 & 6  & 3 \\
D3 --- Purpose       & 2  & 2 & 1  & 0 \\
D4 --- Dependencies  & 8  & 4 & 13 & 0 \\
D5 --- Examples      & 17 & 2 & 17 & 0 \\
D6 --- Contribution  & 9  & 2 & 4  & 1 \\
\hline
\end{tabular}
\end{table}

For Llama, the most pronounced underscoring occurred for D5
(examples and tutorials), with 17 lower scores compared with two
higher scores, and D2 (usage clarity), with 14 lower scores compared
with two higher scores. These dimensions are particularly sensitive
to the README-only evidence policy adopted in this study. Relevant
usage guidance, tutorials, and example material may be provided
through external documentation links that are not followed during
the automated assessment.

Qwen reduced the frequency of disagreement for several dimensions,
particularly D1, D2, D3, and D6. However, systematic underscoring
remained pronounced for D4 and D5. For dependency documentation
(D4), Qwen assigned a lower score than the primary annotator in
13 of 30 repositories and never assigned a higher score. For examples
and tutorials (D5), it assigned a lower score in 17 repositories and
again never assigned a higher score. Against the second annotator,
the D5 tendency was even more pronounced, with Qwen assigning a lower
score in 26 of 30 repositories.

The persistence of the D5 pattern across both model configurations
indicates that examples and tutorial evidence remain particularly
difficult under the implemented README-only assessment design.
The Qwen results additionally identify dependency documentation
(D4) as a recurring source of conservative scoring for the second
model configuration. These patterns show that stronger aggregate
agreement does not eliminate systematic dimension-specific
disagreements.

Purpose description (D3) illustrates a different form of
model-dependent behaviour. Llama produced a negative weighted Kappa
($\kappa=-0.071$) despite 86.7\% exact agreement with the primary
annotator. Human scores for this dimension were strongly concentrated
at score 2, with 28 of 30 repositories receiving that score, making
Kappa particularly sensitive to the small number of disagreements.
Qwen, however, agreed with the primary annotator on 29 of 30
repositories and achieved $\kappa=0.789$, while reaching exact
agreement with the second annotator on all 30 repositories
($\kappa=1.000$). The contrast demonstrates that the weak Llama D3
Kappa was not reproduced by the second model configuration.

Overall, the disagreement patterns reinforce the distinction between
model-level and dimension-level reliability. Qwen achieved stronger
aggregate agreement than Llama in this experiment, but dependencies
and examples remained substantially more difficult than dimensions
such as installation, purpose, and contribution information.

\subsection{Comparison with Inter-Annotator Agreement}

Table~\ref{tab:agreement_summary} places the model--human agreement
results alongside the agreement observed between the two human
annotators.

\begin{table}[htbp]
\centering
\caption{Agreement summary for human--human and model--human
comparisons (macro-average weighted $\kappa$, D1--D6).}
\label{tab:agreement_summary}
\begin{tabular}{lr}
\hline
\textbf{Comparison} & \textbf{Macro-avg. $\kappa$} \\
\hline
Human A1 vs Human A2 & 0.709 \\
Llama vs Human A1    & 0.329 \\
Llama vs Human A2    & 0.275 \\
Qwen vs Human A1     & 0.612 \\
Qwen vs Human A2     & 0.557 \\
\hline
\end{tabular}
\end{table}

The two model configurations show substantially different levels of
agreement with the human annotations. Llama achieved macro-average
weighted Kappa values of 0.329 and 0.275 against the primary and
second annotators, respectively. Qwen achieved higher values of
0.612 and 0.557 against the same two references.

Human--human agreement was 0.709. This value provides useful context
for interpreting the model results, but it should not be treated as
a formal upper bound on achievable model performance. The Qwen
results nevertheless lie considerably closer to the observed
human--human agreement than the Llama results.

The same broad ordering is observed against both human references,
while the dimension-level results vary substantially. This indicates
that reliability cannot be characterised solely as a property of
``LLM-based assessment'' in general. Within the present experiment,
agreement depends materially on the model configuration, the
documentation dimension being assessed, and, to a lesser extent,
the human annotation used as the reference.

\subsection{Model-Reported Confidence Analysis}

Both LLM configurations produced a self-reported confidence label
(HIGH, MEDIUM, or LOW) for each documentation-quality score. To
examine whether these labels provide information about assessment
reliability, the 180 dimension-level predictions produced by each
model (30 repositories $\times$ 6 dimensions) were analysed relative
to the primary human annotations.

For Llama, predictions were distributed across all three confidence
levels. Table~\ref{tab:confidence_calibration} reports exact
agreement, mean absolute error (MAE), and linearly weighted Cohen's
Kappa within each confidence group.

\begin{table}[htbp]
\centering
\caption{Llama agreement with the primary human annotations grouped
by model-reported confidence across all documentation dimensions.}
\label{tab:confidence_calibration}
\begin{tabular}{lrrrr}
\hline
\textbf{Confidence} & \textbf{n} & \textbf{Exact} &
\textbf{MAE} & \textbf{$\kappa$} \\
\hline
HIGH   & 113 & 56.6\% & 0.460 & 0.520 \\
MEDIUM & 60  & 71.7\% & 0.283 & 0.478 \\
LOW    & 7   & 28.6\% & 1.143 & $-$0.120 \\
\hline
\end{tabular}
\end{table}

For Llama, the reported confidence levels were not monotonically
associated with prediction accuracy. MEDIUM-confidence predictions
achieved the highest exact agreement (71.7\%) and lowest MAE (0.283),
whereas HIGH-confidence predictions achieved lower exact agreement
(56.6\%) and a higher MAE (0.460). LOW-confidence predictions
performed worst, with 28.6\% exact agreement and an MAE of 1.143,
although this group contained only seven predictions and should
therefore be interpreted cautiously.

The dimension-level results also show that HIGH confidence did not
consistently indicate stronger agreement. For example, Llama's
HIGH-confidence predictions for D3 achieved 94.4\% exact agreement,
whereas HIGH-confidence predictions for D5 achieved only 29.4\%.
MEDIUM-confidence predictions also exceeded HIGH-confidence
predictions in exact agreement for D1, D2, D4, and D5.

The Qwen configuration exhibited a different confidence pattern.
All 180 dimension-level predictions were assigned HIGH confidence;
no MEDIUM- or LOW-confidence predictions were produced. Consequently,
confidence calibration across confidence levels could not be
evaluated for this model. Across the pooled HIGH-confidence
predictions, exact agreement with the primary annotator was 72.8\%,
MAE was 0.272, and weighted Kappa was 0.662. These pooled values
describe overall performance among predictions that all received the
same confidence label and should not be interpreted as evidence of
confidence calibration.

The uniform Qwen confidence assignment was not reflective of the
substantial variation in observed dimension-level agreement. Exact
agreement ranged from 43.3\% for examples and tutorials (D5) to
96.7\% for purpose description (D3), despite both dimensions
receiving HIGH confidence for every assessed repository. Thus,
Qwen's self-reported confidence did not provide a discriminative
reliability signal in this experiment.

Taken together, the two configurations demonstrate different
limitations of self-reported confidence. Llama produced multiple
confidence levels, but these were not monotonically related to
agreement with human assessment. Qwen produced no confidence
variation at all, preventing discrimination between more and less
reliable predictions. Model-reported confidence is therefore retained
as contextual information in the hybrid FAIRness profile rather than
treated as a calibrated quantitative estimate of assessment
correctness.


\subsection{README Length and Verbosity Bias}

To investigate whether either LLM configuration systematically
assigned relatively higher or lower scores to repositories with
longer README files, README length was compared with the difference
between each model's total documentation score and the corresponding
primary human total. README length was measured as the number of words
in the effective README text supplied to the model after applying the
12,000-character preprocessing limit.

For Llama, Spearman's rank correlation showed a weak negative and
non-significant relationship between effective README length and
model--human score difference ($\rho=-0.169$, $p=0.371$, $n=30$).
For Qwen, the corresponding relationship was similarly weak and
non-significant ($\rho=-0.143$, $p=0.451$, $n=30$). Thus, neither
model showed evidence at the repository level that longer README
files systematically produced relatively higher documentation scores
in the evaluated sample.

For Llama, dimension-level correlations were also non-significant,
with coefficients ranging from $\rho=-0.180$ for dependency
documentation (D4) to $\rho=0.060$ for installation documentation
(D1). In particular, D2 and D5, which showed pronounced Llama
underscoring, had weak negative correlations of $\rho=-0.110$ and
$\rho=-0.135$, respectively.

For Qwen, most dimension-level relationships were likewise weak and
non-significant. D4 was the exception, showing a negative association
between effective README length and model--human score difference
($\rho=-0.392$, $p=0.032$). This direction is consistent with the
systematic Qwen underscoring observed for D4. However, this result is
treated as exploratory because it arises from a dimension-level
analysis across six dimensions with only 30 repositories and no
multiple-comparison correction. It therefore does not provide a
basis for concluding that README length causes the D4 disagreement.

Only one of the 30 README files exceeded the 12,000-character input
limit. Inspection of the excluded portion showed that it consisted
of reStructuredText definitions for images, badges, logos, and
associated links rather than substantive documentation corresponding
to D1--D6. No relevant assessment evidence was identified as lost
through truncation in this particular case. Nevertheless, this
observation concerns only one repository, and the fixed sequential
cutoff may remove relevant sections in other long READMEs.

Overall, the repository-level analyses provide no evidence of a
systematic README verbosity bias for either evaluated model
configuration. This result is limited to the present 30-repository
sample and should not be interpreted as evidence that README length
or truncation effects cannot occur in a broader repository
population.

\section{Hybrid FAIRness Profile Analysis}

The hybrid FAIRness profile presents the rule-based indicators
(F1--R2) and LLM-based documentation quality scores (D1--D6)
together in a single structured record per repository while
preserving them as separate assessment dimensions. No overall
scalar FAIRness score is computed across the two components.
This section analyses the profiles for the 30 ground truth
repositories and examines the relationship between the two
assessment dimensions.
\subsection{Score Distributions}

Table~\ref{tab:hybrid_summary} presents descriptive scores for the
rule-based component and the primary Llama-based documentation
component across the 30 ground truth repositories.

\begin{table}[h]
\centering
\caption{Mean rule-based and primary Llama-based documentation scores
across the 30 ground truth repositories.}
\label{tab:hybrid_summary}
\begin{tabular}{lrr}
\hline
\textbf{Component} & \textbf{Mean Score} & \textbf{Maximum} \\
\hline
Rule-based total & 4.30 & 6 \\
Llama documentation total & 6.70 & 12 \\
\hline
\end{tabular}
\end{table}

The mean rule-based score across the 30 ground truth repositories was
4.30 out of 6, slightly higher than the mean of 3.86 across the full
188-repository dataset. The mean documentation score produced by the
primary Llama configuration was 6.70 out of 12.

These totals are reported only as descriptive summaries within their
respective assessment components. They are not combined or interpreted
as a single overall FAIRness score. Although a corresponding hybrid profile set was also generated using the
Qwen results, the purpose of the second model evaluation is to assess the
robustness of the LLM-based assessment findings rather than to replace the
primary Llama-based profile.

\subsection{Relationship Between Rule-Based and LLM Scores}

To examine whether repositories that score highly on structural FAIR
indicators also score highly on documentation quality, Spearman's rank
correlation was computed between the rule-based total and the documentation
total produced by each LLM configuration across the 30 ground truth
repositories.

For Llama, the correlation was weak and not statistically significant
($\rho=0.175$, $p=0.356$). Applying the same analysis to the Qwen-based
hybrid profiles also yielded a weak, non-significant correlation
($\rho=0.134$, $p=0.480$). No clear monotonic association was therefore
detected between the structural and documentation components for either
model in the evaluated sample.

These results are consistent with the design decision to preserve structural
FAIR-related indicators and qualitative documentation assessment as
separate dimensions of the hybrid profile. A repository may achieve a
high rule-based score through the presence of formal metadata and
repository artefacts while providing limited inline documentation.
Conversely, a repository may provide clear and useful README
documentation while lacking persistent identifiers, structured
metadata, or explicit citation information.

The Qwen analysis uses the same structural scores and the same
30 repositories as the Llama analysis, changing only the model-specific
documentation totals. The similar findings indicate that the absence
of a clear association in this sample is not specific to the primary
Llama configuration. However, non-significant correlations in this
small sample do not establish independence between the two assessment
components.

\subsection{Complementary Assessment Dimensions}

The evaluated repositories exhibited complementary assessment 
patterns. Repositories with high rule-based scores but comparatively low
documentation scores may possess strong structural FAIR-related
evidence while providing limited inline documentation. Conversely, 
other repositories achieve higher documentation scores despite 
lacking formal metadata such as \texttt{CITATION.cff} files or 
persistent identifiers.

By presenting both dimensions in parallel, the hybrid profile 
allows users to identify specific FAIRness gaps and prioritise 
improvements accordingly, rather than receiving a single aggregated 
score that obscures the nature of any deficiencies.

\subsection{Profile Interpretation}

The hybrid FAIRness profile provides information from both structural
and documentation-based assessment that is not represented by either
component alone. A repository with high scores in both components
shows strong coverage of the assessed structural FAIR-related
indicators together with strong README documentation according to the
qualitative assessment rubric.

Conversely, a repository with a high rule-based score but a low
documentation score may satisfy many of the assessed structural
indicators while providing limited documentation for practical reuse.
Similarly, a repository with lower rule-based scores but higher
documentation scores may provide strong README documentation while
lacking structural evidence such as persistent identifiers or
machine-readable citation metadata.

The evaluation with two LLM configurations additionally shows that
the documentation component is sensitive to the model used to produce
it. The hybrid representation itself is model-independent in the
sense that the structural and documentation dimensions remain
separate, but the reliability analysis demonstrates that individual
D1--D6 scores should be interpreted with awareness of the underlying
model configuration and its dimension-specific agreement behaviour.

\section{Baseline Comparison}

This section compares the proposed hybrid assessment framework 
with two existing FAIR assessment approaches. The software-oriented 
howfairis tool is evaluated across all 30 ground truth repositories 
and provides a direct comparison of assessment scope. F-UJI is 
examined separately through a five-repository pilot to determine 
whether a data-oriented FAIR assessment tool can provide a 
meaningful quantitative baseline when applied directly to research 
software repositories hosted on GitHub.

\subsection{F-UJI Baseline Applicability}
\label{sec:fuji_baseline}

Before using F-UJI as a quantitative baseline, its applicability 
to research software repositories was investigated. F-UJI is 
designed primarily for assessing the FAIRness of research data 
objects. The pilot therefore examined whether it could meaningfully 
differentiate research software repositories when provided with 
plain GitHub repository URLs.

Five ground truth repositories with different rule-based and 
documentation quality profiles were submitted directly to F-UJI 
(version 3.5.1, metrics v0.8): \texttt{astropy/astropy} and 
\texttt{sunpy/sunpy} with rule-based scores of 6/6, 
\texttt{gonum/gonum} with a rule-based score of 2/6 and an LLM 
documentation score of 4/12, \texttt{yt-project/yt} with scores 
of 3/6 and 10/12, respectively, and \texttt{scipy/scipy} with 
scores of 5/6 and 6/12. Table~\ref{tab:fuji_pilot} presents the 
results.

\begin{table}[h]
\centering
\caption{F-UJI assessment results compared with the proposed 
framework's rule-based and LLM documentation scores (n=5).}
\label{tab:fuji_pilot}
\begin{tabular}{lrrrrrrr}
\hline
\textbf{Repository} & \textbf{Rule} & \textbf{LLM} & 
\textbf{F-UJI} & \textbf{F} & \textbf{A} & 
\textbf{I} & \textbf{R} \\
\hline
\texttt{astropy/astropy} & 6/6 & 4/12 & 42.31\% & 42.86\% & 42.86\% & 33.33\% & 50.00\% \\
\texttt{sunpy/sunpy}     & 6/6 & 6/12 & 42.31\% & 42.86\% & 42.86\% & 33.33\% & 50.00\% \\
\texttt{gonum/gonum}     & 2/6 & 4/12 & 42.31\% & 42.86\% & 42.86\% & 33.33\% & 50.00\% \\
\texttt{yt-project/yt}   & 3/6 & 10/12 & 42.31\% & 42.86\% & 42.86\% & 33.33\% & 50.00\% \\
\texttt{scipy/scipy}     & 5/6 & 6/12 & 42.31\% & 42.86\% & 42.86\% & 33.33\% & 50.00\% \\
\hline
\end{tabular}
\end{table}

Despite rule-based scores ranging from 2/6 to 6/6 and LLM 
documentation scores ranging from 4/12 to 10/12, all five 
repositories received the same F-UJI score of 42.31\% (11/26) 
and the same FAIR category breakdown. F-UJI also reported that 
the evaluated resources did not identify themselves as datasets, 
with the detected resource types including 
\texttt{softwaresourcecode} and \texttt{object}.


Within this pilot, F-UJI therefore provided no differentiation 
between repositories with substantially different software-specific 
FAIRness profiles. Inspection of the metric-level outputs suggests 
that several F-UJI metrics rely on metadata representations that 
are not directly aligned with repository-level software evidence. 
For example, the license metric for \texttt{astropy/astropy} 
failed despite the repository explicitly identifying a BSD 
3-clause license.

Given the identical results across the five varied repositories 
and F-UJI's resource-type warning, a quantitative comparison 
across all 30 ground truth repositories was not pursued. The pilot 
is instead treated as an applicability finding: when applied
directly to the tested GitHub repository URLs, F-UJI did not provide
sufficient differentiation within this pilot to serve as a quantitative
baseline for the software-oriented assessment performed in this study.

\subsection{howfairis Assessment Results}

The howfairis tool was applied to all 30 ground truth repositories. 
howfairis evaluates five binary indicators: open repository, 
license, software registry, citation, and checklist (presence of 
a fair-software.eu badge). Table~\ref{tab:howfairis} summarises 
the score distribution.

\begin{table}[h]
\centering
\caption{howfairis score distribution across the 30 ground truth 
repositories.}
\label{tab:howfairis}
\begin{tabular}{rrr}
\hline
\textbf{Score} & \textbf{Count} & \textbf{Percentage} \\
\hline
5/5 & 1  & 3\%  \\
4/5 & 15 & 50\% \\
3/5 & 11 & 37\% \\
2/5 & 3  & 10\% \\
\hline
\end{tabular}
\end{table}

Only one repository (\texttt{Ouranosinc/xclim}) achieved a full 
score of 5/5. The checklist indicator, which requires the presence 
of a fair-software.eu badge in the repository README, was satisfied 
by only one repository in the ground truth set. Consequently, this 
indicator provided little differentiation between repositories in 
the evaluated sample. Both the open repository and license 
indicators passed for all or nearly all repositories, consistent 
with the rule-based results reported in Section~5.3.

\subsection{Scope Comparison}

Table~\ref{tab:scope} compares the assessment dimensions covered 
by howfairis, the rule-based component, and the LLM-based component 
of the proposed framework. F-UJI is not included in this 
dimension-by-dimension table because the pilot demonstrated that 
its data-oriented metrics were not directly comparable to the 
software-specific dimensions evaluated by the proposed framework.

\begin{table}[h]
\centering
\caption{Comparison of assessment scope across howfairis, the 
rule-based component, and the LLM-based component.}
\label{tab:scope}
\begin{tabular}{lccc}
\hline
\textbf{Dimension} & \textbf{howfairis} & \textbf{Rule-based} & 
\textbf{LLM-based} \\
\hline
Open repository       & \checkmark & $-$ & $-$ \\
License               & \checkmark & \checkmark (A1) & $-$ \\
Software registry     & \checkmark & $-$ & $-$ \\
Citation metadata     & \checkmark & \checkmark (F2, R2) & $-$ \\
Checklist badge       & \checkmark & $-$ & $-$ \\
Persistent identifier & $-$ & \checkmark (F1) & $-$ \\
Dependency file       & $-$ & \checkmark (I1) & $-$ \\
Installation clarity  & $-$ & $-$ & \checkmark (D1) \\
Usage clarity         & $-$ & $-$ & \checkmark (D2) \\
Purpose description   & $-$ & $-$ & \checkmark (D3) \\
Dependency description & $-$ & $-$ & \checkmark (D4) \\
Examples/tutorials    & $-$ & $-$ & \checkmark (D5) \\
Contribution info     & $-$ & $-$ & \checkmark (D6) \\
\hline
\end{tabular}
\end{table}

The comparison shows that howfairis and the proposed framework 
cover largely complementary assessment dimensions. Both approaches 
assess licensing and citation-related metadata, which represents 
the primary overlap between them. The proposed framework 
additionally evaluates persistent identifiers and dependency 
specifications through its rule-based component, as well as six 
documentation quality dimensions through its LLM-based component. 
These aspects are not assessed by howfairis.

During the evaluation, the software registry indicator appeared to 
depend on repository badges displayed in the README rather than 
direct verification of registry entries. As a result, repositories 
available through a software registry but lacking the corresponding 
badge received a negative registry score.

\section{Summary of Findings}

The evaluation produced five main findings.

First, the rule-based component showed complete agreement with the
finalised primary annotations for all six structural indicators in
the 30-repository ground truth corpus. Precision, recall, and F1-score
were 1.00 for F1, F2, A1, I1, R1, and R2. The same result was observed
on the 20-repository pilot-exclusion subset. Because some structural
annotations in this subset were corrected during final scorer
validation, this result is interpreted as an internal consistency
check rather than independent validation evidence.

Second, the reliability of qualitative documentation assessment was
strongly dependent on both the LLM configuration and the documentation
dimension. Against the primary annotator, Llama achieved a
macro-average weighted Cohen's Kappa of 0.329, whereas Qwen achieved
0.612. Against the second annotator, the corresponding values were
0.275 and 0.557. Human--human agreement across D1--D6 was 0.709,
providing contextual evidence that the two human annotators themselves
did not agree perfectly on all qualitative dimensions.

The dimension-level results further demonstrate this model
dependence. Qwen showed particularly strong agreement with the primary
annotator for installation (D1, $\kappa=0.799$), purpose
(D3, $\kappa=0.789$), and contribution information
(D6, $\kappa=0.759$). The contrast for purpose was especially
pronounced: Llama achieved $\kappa=-0.071$, whereas Qwen achieved
$\kappa=0.789$. Thus, the weak Llama result for this dimension was
not reproduced by the second model configuration.

At the same time, stronger aggregate Qwen agreement did not eliminate
dimension-specific difficulties. Dependency documentation (D4) and
examples or tutorials (D5) remained more difficult, with Qwen
achieving $\kappa=0.487$ and $\kappa=0.320$, respectively, against
the primary annotator. Qwen assigned lower D4 scores than the primary
annotator in 13 of 30 repositories and lower D5 scores in 17 of 30,
without assigning a higher score in either dimension. Llama also
showed pronounced underscoring for D5. These results indicate that
qualitative LLM assessment cannot be characterised by a single
reliability value; performance varies materially across both model
configurations and documentation dimensions.

Third, the pilot-exclusion sensitivity analysis showed that the
model--human agreement findings were not primarily driven by the
10 repositories used during pilot-phase refinement. Llama's
macro-average weighted Kappa increased from 0.329 on the complete
30-repository corpus to 0.373 on the 20-repository subset. Qwen
remained almost unchanged, from 0.612 to 0.615. The latter result
provides particularly clear evidence that the stronger Qwen agreement
was stable to exclusion of the original pilot repositories.

Fourth, model-reported confidence did not provide a reliable
discriminative signal of assessment correctness. For Llama,
MEDIUM-confidence predictions achieved higher exact agreement than
HIGH-confidence predictions, showing that confidence was not
monotonically related to observed accuracy. Qwen assigned HIGH
confidence to all 180 dimension-level assessments and produced no
MEDIUM- or LOW-confidence predictions. Consequently, calibration
across confidence levels could not be evaluated for Qwen. Despite
the uniform HIGH-confidence assignment, Qwen exact agreement varied
from 43.3\% for D5 to 96.7\% for D3. Self-reported confidence should
therefore not be interpreted as a calibrated probability of
correctness in the present framework.

Relatedly, no repository-level evidence of systematic README
verbosity bias was observed for either model configuration.
The correlation between effective README length and model--human score difference was weak
and non-significant for Llama ($\rho=-0.169$, $p=0.371$) and Qwen
($\rho=-0.143$, $p=0.451$). An exploratory Qwen dimension-level
analysis produced a negative association for D4
($\rho=-0.392$, $p=0.032$), but this result is not treated as
confirmatory because of the small sample and multiple
dimension-level comparisons.

Fifth, the hybrid profile captures complementary aspects of research
software assessment. The correlation between the rule-based total and
documentation total was weak and non-significant for both Llama
($\rho=0.175$, $p=0.356$) and Qwen ($\rho=0.134$, $p=0.480$).
These findings are consistent with structural FAIR-related evidence
and qualitative README documentation providing different information,
although they do not establish statistical independence. The
framework consequently retains them as separate components rather
than collapsing them into a single scalar FAIRness score.

The baseline experiments further support this distinction in scope.
The software-oriented howfairis comparison showed overlap with some
structural indicators while covering a narrower set of repository
properties. The F-UJI pilot demonstrated limited applicability of a
data-oriented FAIR assessment service when applied directly to the
selected software repositories. These comparisons support the need
for an assessment design that is explicitly adapted to research
software while retaining transparent, criterion-level evidence.

Overall, the evaluation answers the two research questions from
complementary perspectives. For RQ1, the LLM-based component adds
qualitative information about installation, usage, purpose,
dependencies, examples, and contribution guidance that is not
captured by deterministic repository checks. The weak relationship
between the structural and documentation components further supports
their complementary role.

For RQ2, the two-model evaluation shows that LLM-based documentation
assessment can achieve substantial agreement with human judgement for
some model--dimension combinations, but reliability is neither uniform
across dimensions nor stable across model configurations. Qwen
achieved substantially stronger overall agreement than the primary
Llama configuration in this experiment, while D4 and D5 remained
comparatively difficult and self-reported confidence did not provide
a dependable reliability signal. Human review and transparent
criterion-level evidence therefore remain important when interpreting
LLM-generated FAIRness assessments.
